% ACL body: condensed from ../latex/body.tex around the main argument.
\section{Introduction}

A model can return a valid, well-typed decision and still be wrong. When
several such decisions are chained into a workflow, errors propagate, and
a common hope is that the steps' reported confidence can say when to hand
a case to a human. Typed decision models make this hope concrete. Jev,
TypeSafe's decision API, answers closed questions through three
primitives: a categorical \choice{}, a binary \noul{} and an ordinal
\score{}, each with a number between 0 and 1. Prior work has benchmarked
these primitives and their calibration on single requests
\citep{deusser2026jev,meng2026jevaudit,denney2026jev}, and argued in
general that error risk is not review value \citep{park2026riskreview}.

We study an operational question that single-request benchmarks leave
open. \emph{When a workflow chains heterogeneous typed decisions, which
signals identify the cases where it fails, and can those signals identify
the specific costly errors for which human review is worthwhile?} The two
halves can come apart: a signal can rank failures well and still not say
which reviews pay.

We chain Jev's three primitives into a dependent bug-triage workflow
(route, assess regression, rate severity) and study it on 202 Mozilla
bugs labeled by triage and 240 Linux kernel bugs labeled by their
reporters, followed by pre-registered tests on 500 new bugs and four
further models. Our contributions:
\begin{enumerate}
  \item \textbf{Workflow failure prediction} (\cref{sec:heldout,sec:crossmodel}).
  A frozen decision-feature estimator ranks wrong workflow actions better
  than final-step confidence on 500 new bugs, and the gap replicates
  across five models from four sources. Most of the signal is the chosen
  action's historical error rate.
  \item \textbf{Prediction is not review value} (\cref{sec:intervention}).
  Under the pre-registered costs, no review strategy built on that signal
  beats no review at a 10\% budget. An oracle locates the value in wrongly
  accepted cases that the tested signals do not isolate, and an
  exploratory sensitivity analysis shows that whether review pays at all
  depends on the ratio of wrong-accept cost to review cost. A case-level
  score failed a pre-registered feasibility check, so its confirmatory
  study was stopped by rule before new data.
  \item \textbf{Supporting evidence} (\cref{sec:primitives}). Inside the
  workflow the three primitives' confidence signals behave differently,
  thresholds set on constructed data do not transfer, and \score{}
  compresses severity on both trackers. These results agree with the
  single-request audits and help explain the workflow-level findings.
\end{enumerate}
All data collection and confirmatory analyses were pre-registered: each
protocol was committed, with its SHA-256 hash, before the data it governs
existed. The anonymized supplementary archive lists the commit dates and
hashes, and deviations are disclosed (\cref{app:deviations}).

\section{Related work}

\para{Evaluations of Jev and open decision models} Jev has been
benchmarked zero-shot on 37 datasets, where its \choice{} probabilities
were well calibrated and its binary probabilities ranked well but sat
poorly relative to a 0.5 threshold \citep{deusser2026jev}. A calibration
audit (described by its author as exploratory) found that calibration
changed with the question form, that 71.6\% of four-class answers sat at
exactly 1.00, leaving a top-class threshold only 53 operating points, and
that no threshold certified a 5\% automatic tier on the four-class task
\citep{meng2026jevaudit}. On political-science replications Jev's
calibration advantage was not consistent \citep{denney2026jev}, and an
early evidence audit found no independent accuracy advantage yet for the
typed readout \citep{tang2026typed}. Open instruction-tuned models give
Jev-like decisions by reading label logits in one forward pass
\citep{geuter2026innerjev}, the readout we use for the general models.
These studies evaluate single requests; none studies \score{} inside a
workflow, composition of dependent decisions, or the cost of human review.

\para{Review value and deferral} Learning to defer
\citep{madras2018predict}, selective prediction
\citep{geifman2017selective}, cost-sensitive deferral under workload
limits \citep{alves2024defer} and cascades with escalation
\citep{jung2025trust,bouchard2026escalation} address deferral of single
decisions. Closest to our intervention result, \citet{park2026riskreview}
argue that review priority should depend on wrongness, repairability,
impact and cost rather than risk alone, and show on a QA benchmark that
review-value ranking leaves less residual error than risk ranking. Our
result is an instance of that distinction in a dependent typed-decision
workflow, with a frozen cost table, pre-registered tests and a
feasibility check for a case-level policy.

\para{Composition and trajectories} Errors compound across pipeline
stages, and independently calibrated stages lack joint guarantees
\citep{kotte2026pasc,chowdhury2019error}; step-level calibration need not
imply trajectory-level calibration for LLM agents \citep{lin2026uqagents};
and the product rule for combining confidences assumes conditional
independence \citep{kittler1998combining}. LLM raters compress ordinal
scores toward the middle \citep{zhang2026centraltendency}.

\para{Positioning} We are not aware of prior work that tests, in a
dependent workflow of structurally different typed decisions with real
labels, held-out data and several models, which signals predict workflow
failure and whether they identify the reviews worth paying for. We make
no priority claim: our search was web-based, and the cited Jev studies
appeared in September and October 2026.

\section{Setting}
\label{sec:setting}

\begin{figure}[t]
\centering
\resizebox{\linewidth}{!}{%
\begin{tikzpicture}[node distance=0.55cm, every node/.style={font=\small},
  box/.style={draw, rounded corners, minimum width=1.6cm, minimum height=0.9cm, align=center},
  arr/.style={-{Stealth[length=2mm]}, thick}]
\node[box, fill=blue!6] (ev) {bug\\ text $x$};
\node[box, right=of ev] (c) {\choice{}\\ route $c_1$};
\node[box, right=of c] (n) {\noul{}\\ regression $c_2$};
\node[box, right=of n] (s) {\score{}\\ severity $c_3$};
\node[box, below=0.6cm of n, fill=orange!8] (f) {action $a = F(c_2, c_3)$};
\draw[arr] (ev) -- (c);
\draw[arr] (c) -- (n);
\draw[arr] (n) -- (s);
\draw[arr] (n) -- (f);
\draw[arr] (s) |- (f);
\end{tikzpicture}}
\caption{The workflow. Each node reads the bug text plus all upstream
decisions and is called three times. The action rule $F$ uses $c_2$ and
$c_3$; a review replaces the action with \emph{escalate}.}
\label{fig:workflow}
\end{figure}

\para{Workflow} Node~1 (\choice{}) routes a bug to the network or storage
subsystem; node~2 (\noul{}) assesses whether it is a regression; node~3
(\score{}) rates severity on three levels. Each node reads the bug text
plus the upstream decisions, appended by a fixed adapter
(\cref{fig:workflow}), and is called three times; the answers are reduced to the modal decision (mean
probability for \noul{}). All calls used \texttt{jev-1.13.0}.

\para{Target action and cost} The action rule $F(c_2, c_3)$ maps an
unclear regression answer to \emph{request evidence}, severity 2 to
\emph{escalate}, severity 1 with a regression to \emph{defer}, and
everything else to \emph{accept} (\cref{tab:F}). The target is
$\astar = F$ applied to the reference labels. It is a construct we
defined, and every strategy result is conditional on it. Costs come from a
frozen matrix (\cref{tab:cost}): a wrong \emph{accept} costs 6--10, any
other mismatch 1, and each escalation an extra 1 for human time. A review
turns a case's action into \emph{escalate}, so reviewing a wrongly
accepted case saves 6 or 9, and any other review loses 1--2.

\para{Data} \emph{Mozilla:} 202 fixed Firefox networking and storage
bugs; severity is assigned by triage and the regression flag is a
keyword. \emph{Linux kernel:} 240 fixed networking and storage bugs;
severity is set by the reporter and the regression flag is a field. Both
development samples are balanced across (subsystem, level) cells. The
pre-registered tests use 500 \emph{new} bugs, 250 per tracker, sampled at
each tracker's natural mix. Agreement is with reference labels, not
ground truth.

\para{Statistics} AUROCs use the Mann--Whitney estimator with percentile
bootstrap intervals over bugs ($B = 2{,}000$, stratified by tracker).
Confirmatory tests are one-sided paired bootstrap tests with Holm
correction within each pre-registered family.

\section{Supporting evidence: the primitives inside the workflow}
\label{sec:primitives}

\begin{table}[t]
\centering
\small
\setlength{\tabcolsep}{2pt}
\begin{tabular}{lcc}
\toprule
 & Mozilla (triage) & Kernel (reporter) \\
\midrule
\choice{} at $\geq 0.90$ & 85\% & 96\% \\
\choice{} conf. AUROC & 0.90 [0.83, 0.96] & 0.98 [0.95, 1.00] \\
\noul{} value AUROC & 0.53 [0.42, 0.65] & 0.41 [0.32, 0.50] \\
\score{} conf.\ AUROC & 0.58 [0.50, 0.66] & 0.57 [0.50, 0.64] \\
\midrule
Agreement: \choice{} & 0.89 [0.83, 0.92] & 0.975 [0.95, 0.99] \\
\phantom{Agreement:} \noul{} & 0.80 [0.74, 0.85] & 0.725 [0.67, 0.78] \\
\phantom{Agreement:} \score{} & 0.53 [0.46, 0.60] & 0.44 [0.38, 0.50] \\
\bottomrule
\end{tabular}
\caption{Each primitive's confidence signal, as AUROC for predicting the
node's own agreement with the reference label (95\% bootstrap CI), and
node agreement. Full table and reliability diagrams in \cref{app:primitives}.}
\label{tab:prim}
\end{table}

The three primitives' signals are not one kind of quantity
(\cref{tab:prim}). \choice{} confidence is saturated (85--96\% of bugs at
$\geq 0.90$) yet discriminative, because its rare low-confidence answers
are where its errors are. \noul{}'s confidence-like value $2|\bar p - 0.5|$
carries no usable information about its own correctness. \score{}'s
confidence is flat across true severity levels, so it gives no warning of
a systematic bias (\cref{fig:sev}): on both trackers most minor bugs are rated higher than
their label (79\% [0.67, 0.88] and 75\% [0.65, 0.83]) and many serious bugs lower (43\% and 56\%), a
central-tendency pattern reported for LLM raters elsewhere
\citep{zhang2026centraltendency}. These findings agree with the
single-request audits where they overlap: \choice{} is saturated yet
discriminative, and thresholds must be validated on target data. On
\noul{}, the audits report accurate, recalibratable yes/no answers on
sentiment, whereas on our regression question its probability did not
discriminate its own correctness, so binary-probability quality looks
task-dependent.

\begin{figure}[t]
\centering
\begin{tikzpicture}
\begin{groupplot}[group style={group size=2 by 1, horizontal sep=0.35cm, y descriptions at=edge left},
  width=0.5\linewidth, height=4.2cm, view={0}{90}, enlargelimits=false, axis on top,
  colormap={sev}{rgb255=(255,255,255) rgb255=(107,174,214) rgb255=(33,113,181)},
  point meta min=0, point meta max=0.8,
  xtick={0,1,2}, ytick={0,1,2}, y dir=reverse,
  yticklabels={minor,normal,serious}, xticklabels={0,1,2},
  tick label style={font=\scriptsize}, label style={font=\scriptsize}, title style={font=\small},
  xlabel={\score{} output},
  nodes near coords={\pgfmathprintnumber[fixed,fixed zerofill,precision=2]{\pgfplotspointmeta}},
  nodes near coords style={font=\scriptsize, anchor=center}]
\nextgroupplot[title={Mozilla (triage)}, ylabel={true level}]
\addplot[matrix plot*, mesh/cols=3, point meta=explicit] table[meta=v] {fig_heat_mozilla.dat};
\nextgroupplot[title={Kernel (reporter)}]
\addplot[matrix plot*, mesh/cols=3, point meta=explicit] table[meta=v] {fig_heat_kernel.dat};
\end{groupplot}
\end{tikzpicture}
\caption{Row-normalized confusion of \score{} against reference severity
(repeat-unstable outputs excluded). On both trackers the middle column
absorbs mass from both extremes: minor bugs are rated up, serious bugs
down.}
\label{fig:sev}
\end{figure}

\begin{figure}[t]
\centering
\begin{tikzpicture}
\begin{axis}[width=0.78\linewidth, height=5.0cm, xmin=0.25, xmax=1.05, ymin=0.4, ymax=5.6,
  ytick={1,2,3,4,5}, yticklabels={(B) any error, (A) missed serious, \score{} own, \noul{} own, \choice{} own},
  xlabel={AUROC (95\% bootstrap interval)}, grid=major, grid style={gray!20},
  tick label style={font=\scriptsize}, label style={font=\scriptsize},
  legend style={font=\scriptsize, at={(0.98,0.02)}, anchor=south east}]
\draw[dashed, gray] (axis cs:0.5,0.4) -- (axis cs:0.5,5.6);
\addplot[blue, only marks, mark=*, error bars/.cd, x dir=both, x explicit]
  table[x=x, y=y, x error minus=lo, x error plus=hi] {fig_forest_mozilla.dat};
\addplot[orange, only marks, mark=square*, error bars/.cd, x dir=both, x explicit]
  table[x=x, y=y, x error minus=lo, x error plus=hi] {fig_forest_kernel.dat};
\legend{Mozilla, Kernel}
\end{axis}
\end{tikzpicture}
\caption{Development-set AUROCs with 95\% intervals; dashed line: chance.
Only \choice{}'s confidence discriminates its own errors. Composed
confidence is at chance for both escalation targets, (A) and (B).}
\label{fig:forest}
\end{figure}

\para{Fixed thresholds do not transfer} Two escalation thresholds designed
on constructed data lost their selectivity on real text: a \choice{}
trigger fired on 85--96\% of bugs everywhere, and a confidence-product cut
that escalated 73\% of constructed items (target 67\%) escalated 88\% and
78\% of real bugs (targets 41--42\%), because real text pulls \noul{}
toward 0.5.

\para{Composed confidence does not pick escalations (exploratory)}
Escalating, on top of the decision-based action, the bugs with the lowest
composed confidence did no better than random selection on either tracker.
Composed confidence ranked missed serious bugs at AUROC 0.49 [0.37, 0.61]
and 0.50 [0.40, 0.60], and any decision error at 0.54 [0.45, 0.63] and
0.47 [0.40, 0.54] (\cref{fig:forest}). Only \choice{}'s confidence
discriminates, and it belongs to the one node the action rule does not
use. The obvious comparison, decision-based against
confidence-based composition, is confounded, because the target is itself
a function of the decisions. These analyses were specified after the
Mozilla results and motivated the pre-registered tests below
(\cref{app:exploratory}).

\section{A pre-registered held-out test}
\label{sec:heldout}

We trained one decision-feature estimator ($E_{\text{dec}}$:
$L_2$-regularized logistic regression on the \score{} level, the \noul{}
decision and the workflow action) on all 442 development bugs and froze it
by hash before any new data existed. We then collected the 500 new bugs
(4{,}500 Jev calls). Pre-registered hypotheses: HR1, $E_{\text{dec}}$ ranks
wrong workflow actions better than low final-step confidence; HR2, it
ranks above chance on each tracker.

Both are supported. HR1: $E_{\text{dec}}$ reaches AUROC 0.763 against 0.549
for final-step confidence, a difference of $+0.214$ (one-sided 95\% lower
bound $+0.158$). HR2: AUROC is 0.811 on Mozilla (lower bound 0.762) and
0.697 on the kernel (lower bound 0.648). The estimator's predicted failure
rate matches the observed rate overall, but it outputs only 7--8 distinct
risk levels (ECE 0.18--0.20), so it is a ranking, not a calibrated
probability. Held-out and development bugs share no bug IDs, no exact
texts and no near-duplicates (word 5-gram Jaccard $\geq 0.8$; the highest
similarity of any held-out bug to a development bug is 0.58). Tracker
links (duplicates, see-also, blocks, depends-on, regressions) connect 9
held-out bugs to development bugs; without them $E_{\text{dec}}$ reaches
0.770 against 0.551. On the 70 Mozilla bugs filed in the 12 months before
sampling, which are less likely to be in any model's training data,
$E_{\text{dec}}$ reaches 0.847 against 0.519.

\section{Cross-model replication}
\label{sec:crossmodel}

We ran the same 942 bugs, questions and adapter through four further
models. \textbf{OpenJev 27B} \citep{openjev2026} is an independently
trained Jev-style decision model (Hugging Face \texttt{openjev/openjev},
commit \texttt{b4a89f5}; unrelated projects share the name), served on
one H100 exactly as its authors specify. \textbf{Qwen2.5-7B}
\citep{qwen2024qwen25}, \textbf{OLMo-2-7B} \citep{olmo2024olmo2} and
\textbf{Phi-3.5-mini} \citep{abdin2024phi3} run behind a Jev-compatible
readout: next-token probabilities over the allowed labels, with
TypeSafe's documented confidence formulas. Each model's estimator was
fitted on its own outputs for the 442 development bugs and tested on the
500 held-out bugs; the tests mirror HR1 and HR2, Holm-corrected across
all 8. Three amendments are disclosed in \cref{app:deviations}, including
a readout bug fixed after interim results, with all general models re-run.

\begin{table*}[t]
\centering
\resizebox{\textwidth}{!}{\input{table_crossmodel.tex}}
\caption{Per-model results on the same 500 held-out bugs. AUROC columns
rank wrong workflow actions (95\% bootstrap CI): the decision-feature
estimator fitted on each model's own development outputs, the model's
final-step confidence, the action-conditional error-rate rule, and Jev's
frozen estimator applied unchanged. Cost columns are from the
pre-registered intervention test (\cref{sec:intervention}); the oracle is
an upper bound that must spend all 50 reviews, and ``wrong accepts'' is
their share of no-review cost. Absolute failure rates are not a clean
quality comparison, because the Jev-style and general models use
different interfaces.}
\label{tab:crossmodel}
\end{table*}

\begin{figure}[t]
\centering
\begin{tikzpicture}
\begin{axis}[width=0.8\linewidth, height=4.4cm, xmin=0.3, xmax=0.9, ymin=0.4, ymax=5.6,
  ytick={1,2,3,4,5}, yticklabels={OLMo-2, Phi-3.5, Qwen2.5, OpenJev, Jev 1.13},
  xlabel={held-out AUROC for wrong actions (95\% CI)}, grid=major, grid style={gray!20},
  tick label style={font=\scriptsize}, label style={font=\scriptsize},
  legend style={font=\scriptsize, at={(0.02,0.02)}, anchor=south west}]
\draw[dashed, gray] (axis cs:0.5,0.4) -- (axis cs:0.5,5.6);
\addplot[blue, only marks, mark=*, error bars/.cd, x dir=both, x explicit]
  table[x=dec, y expr=\thisrow{y}+0.12, x error minus=declo, x error plus=dechi] {fig_crossmodel.dat};
\addplot[orange, only marks, mark=square*, error bars/.cd, x dir=both, x explicit]
  table[x=conf, y expr=\thisrow{y}-0.12, x error minus=conflo, x error plus=confhi] {fig_crossmodel.dat};
\legend{decision estimator, final-step confidence}
\end{axis}
\end{tikzpicture}
\caption{Across five models, the decision-feature estimator ranks wrong
workflow actions better than the model's own final-step confidence. The
intervals overlap only for OLMo-2.}
\label{fig:crossmodel}
\end{figure}

All 8 pre-registered tests are supported (\cref{fig:crossmodel,tab:crossmodel}).
Decision estimators reach AUROC 0.64--0.76 and final-step confidence
0.40--0.58; the weakest model is OLMo-2 (0.637 vs.\ 0.551; one-sided
$p = 0.037$, passing at the last Holm step). Jev's frozen estimator, never
refitted, reaches 0.76 [0.72, 0.80] on OpenJev's outputs, as well as
OpenJev's own estimator does. The $E_{\text{dec}}$ intervals of Jev,
OpenJev and Qwen2.5 overlap, so the models cannot be ranked against each
other; only the within-model contrast is the finding.

\para{The signal is simple} A class-frequency rule, the historical error
rate of the action the workflow chose, estimated on development bugs only,
matches the estimator for Jev (0.76) and OpenJev (0.75) and comes close
for the others (0.58--0.70). The predictive information is mostly
\emph{which action was chosen} and how often that action has been wrong.
Part of this is structural: the action and the target are both outputs of
the same rule $F$, so an action's error rate is a property of the
workflow as much as of the model. We therefore do not read the result as
the estimator being sophisticated, but as evidence that per-step
confidence misses information the workflow itself already exposes.
Paired bootstrap contrasts make this precise. For Jev and OpenJev the
decision estimator adds nothing detectable beyond the action rule
($+0.000$ [$-0.012$, $+0.011$] and $+0.002$ [$-0.007$, $+0.011$]); for
Qwen2.5 and OLMo-2 it adds a little ($+0.035$ and $+0.057$, intervals
excluding 0). Adding the uncertainty features to the decision features
raises held-out AUROC by +0.006 (Jev), +0.004 (OpenJev) and +0.012
(Qwen2.5), none distinguishable from 0, but by +0.054 (Phi-3.5) and +0.078
(OLMo-2), both with intervals excluding 0. For the Jev-style models
confidence is redundant with the decisions; for the weaker general models
raw token probabilities carry some extra information.

\begin{table}[t]
\centering
\small
\resizebox{\linewidth}{!}{\input{table_target_robustness.tex}}
\caption{Exploratory: held-out AUROC of the decision-only estimator on two
targets that do not pass through the action rule $F$, fitted on the
quota-sampled development bugs (``dev fit'') or on the same bugs
re-weighted to the natural label mix, against final-step confidence. Bold:
re-weighted estimator above confidence.}
\label{tab:targets}
\end{table}

\para{Beyond the constructed target (exploratory)} Because the target
$\astar$ is a construct, we also asked whether decisions predict failures
on two targets that bypass $F$: the severity answer being wrong, and any
node being wrong (\cref{tab:targets}). Fitted on the development bugs,
several decision estimators rank held-out severity errors far \emph{below}
chance (0.28 for Jev and OpenJev). The cause is label-prior shift: a
decision's error rate depends on the label mix. In the quota-sampled
development set Jev's level-1 answers are 60\% wrong; in the natural-mix
held-out set, where most bugs are level 1, they are 20\% wrong. Re-weighting
the development bugs to the natural label mix repairs this: the decision
estimator then beats final-step confidence on both targets for four of
five models (Jev 0.79 vs.\ 0.61 on severity), the exception being OLMo-2
on any node. Decision-conditioned error rates are therefore informative
beyond the constructed target, but only when estimated under the
deployment label distribution.

\para{Failure modes differ by model family} Jev and OpenJev both pull
severity to the middle level. The general instruct models push severity
to ``blocking'' for 67--100\% of bugs whatever the true level, and are
wrong far more often (72--80\% vs.\ 43--47\%).

\section{Does prediction make review pay?}
\label{sec:intervention}

\para{Design} Before computing any cost outcome, we pre-registered a
review test on the held-out bugs: 50 reviews (10\%) per model, each
escalating a case to a human and priced by the frozen matrix. The primary
arm (A1) ranks candidates by the \emph{action-level expected saving} of a
review, estimated on development bugs; it is cost-aware rather than
error-rate-based, because reviewing a wrong \emph{defer} never saves
anything. Comparators: confidence-ranked review (B) and random review (C);
Holm across 10 tests.

\begin{table}[t]
\centering
\small
\resizebox{\linewidth}{!}{\input{table_intervention_tests.tex}}
\caption{Pre-registered intervention tests: difference in total cost
(A1 minus comparator; negative favors A1) at a 10\% budget, with 95\%
bootstrap intervals. $^{*}$Supported after Holm correction across all
10 tests.}
\label{tab:itests}
\end{table}

\para{Result} Two of ten pre-registered tests are supported (\cref{tab:itests}): A1 beats
random for Jev and OpenJev, by about 10--12 cost units. A1 is not shown to
beat confidence for any model (for Jev its point estimate is 4 units
lower, with an interval spanning zero). Under the frozen table, every
review arm costs more than no review at all; for Jev, 658 with no review against 681--714 with review
(\cref{fig:intervention}).

\begin{table}[t]
\centering
\small
\resizebox{\linewidth}{!}{\input{table_cost_sensitivity.tex}}
\caption{Exploratory cost sensitivity: best feasible arm minus no review
(negative, in bold: review pays), with rankings frozen as pre-registered.
$m$ scales wrong-accept costs, $h$ is the human cost per escalation, and
$p^{*}$ is the break-even precision for reviewing an accepted case whose
target is \emph{escalate}. $^{\dagger}$The pre-registered table. Point
estimates only.}
\label{tab:costsens}
\end{table}

\para{Sensitivity to costs (exploratory)} The pre-registered result holds
under one frozen cost table. In a post-hoc analysis planned before
computing it, we re-scored the same frozen rankings on a grid of cost
ratios within a factor of 2 of that table (\cref{tab:costsens}). The
conclusion is cost-robust only for OLMo-2. For Jev, OpenJev and Qwen2.5,
review pays when wrong accepts cost twice as much (for Jev, 44--69 units
below no review) or when review costs half as much, because the
break-even precision $p^{*}$ then falls to about the 14--16\% share of
wrongly accepted cases. Some reversals are small (Jev $-2$ at $m = 1$,
$h = 0.5$) and lie well inside the intervals of the original tests. The
ordering is more stable than the sign: the action rule beats random for
both Jev-style models in 12 of 16 cells. The stopping decision below was
made under the frozen table and, by its rule, is not reopened.

\begin{figure}[t]
\centering
\begin{tikzpicture}
\begin{axis}[ybar, bar width=2.6pt, width=\linewidth, height=4.2cm,
  symbolic x coords={Jev,OpenJev,Qwen2.5,Phi-3.5,OLMo-2}, xtick=data, ymin=0, ymax=960,
  ylabel={total cost (500 bugs)}, enlarge x limits=0.12,
  grid=major, grid style={gray!20}, tick label style={font=\scriptsize}, label style={font=\scriptsize},
  legend image code/.code={\draw[#1] (0cm,-0.08cm) rectangle (0.18cm,0.12cm);},
  legend style={font=\scriptsize, at={(0.5,1.02)}, anchor=south, legend columns=3, draw=none}]
\addplot[fill=gray!40] table[x=model, y=noreview] {fig_intervention.dat};
\addplot[fill=blue!55] table[x=model, y=A1] {fig_intervention.dat};
\addplot[fill=orange!70] table[x=model, y=B] {fig_intervention.dat};
\addplot[fill=gray!75] table[x=model, y=C] {fig_intervention.dat};
\addplot[fill=green!55!black] table[x=model, y=oracle] {fig_intervention.dat};
\legend{no review, action rule, confidence, random, oracle}
\end{axis}
\end{tikzpicture}
\caption{Total cost at a 10\% review budget under the pre-registered
cost table. No feasible strategy beats no review. A perfect reviewer (oracle, an upper bound using held-out
outcomes) would cut Jev's cost by 44\% and OpenJev's by 33\%; for the
general models, which escalate most cases themselves, it saves little or
nothing.}
\label{fig:intervention}
\end{figure}

\para{Where the cost is} For the Jev-style models, 40--52\% of no-review
cost comes from wrongly \emph{accepted} serious or regression bugs. An
oracle that picked the right 50 cases would save 33--44\%. But those cases
are 15--16\% of accepted cases. The action rule cannot separate them,
because every accept looks alike to it, and confidence cannot either. The
value exists; the identification fails. For the general models, cost is
dominated by over-escalation, which more review cannot reduce. Ranking
which actions are wrong is therefore not the same as finding which
reviews save money: frequent errors and expensive errors differ, the
distinction \citet{park2026riskreview} draw in general.

\begin{table}[t]
\centering
\small
\setlength{\tabcolsep}{4pt}
\resizebox{\linewidth}{!}{\input{table_budget.tex}}
\caption{Pre-registered feasibility check for the case-level score:
simulated power to show lower cost than no review (one-sided
$\alpha = 0.05$), and the median and 80\% range of the projected cost
change, over 400 simulated studies per cell under the natural label mix.
Viability required power $\geq 0.80$; no cell met it.}
\label{tab:budget}
\end{table}

\para{A case-level score, and a stopping decision} We then asked whether a
per-case score could find the expensive errors. Under a plan committed
before any score was computed, a parameter-free plug-in score was selected
on the 442 development bugs over a small logistic model and frozen. It
ranks accepted cases by expected saving from review,
$9p_2 + 6p_1q - 2(1 - p_2 - p_1q)$, using \score{}'s mass on levels 1 and 2
($p_1, p_2$) and \noul{}'s probability $q$. On the development mix it cut
cost by 12\% against no review (630 vs.\ 713) and captured 22\% of the
oracle saving. But the development bugs were quota-sampled; reweighted to
the natural label mix, a 10\% budget was projected to cost \emph{more}
than no review. Before generating any fresh case, we pre-registered a
feasibility check that varied only the budget (1--10\%), with a binding
stop rule: viability required power $\geq 0.80$ to show lower cost than no
review with at most 1{,}000 fresh bugs, and a negative median projected
cost change, with uncertainty in error prevalence and ranking propagated
by resampling. No budget was viable. The highest power was 0.20 (1\%
budget, 1{,}000 bugs), with a median change of $-2$ and an 80\% range of
$[-46, +20]$; every budget of 2\% or more had a positive median cost change
(\cref{tab:budget}). The confirmatory study was therefore not run, no
further designs were searched, and no fresh API calls were made. This is a
feasibility assessment, not evidence that review cannot help: a good
selector would save substantially, but under the pre-registered costs the
signals tested here did not identify the cases well enough to pay for the
reviews.

\section{Discussion}

\para{Relation to prior work} The single-request audits
\citep{deusser2026jev,meng2026jevaudit} establish that Jev's scores must
be validated per task and question form; our primitive-level results add
that, once chained, the only discriminative confidence can sit on a node
the action never reads, and that \score{}'s confidence is silent about its
own systematic bias. \citet{park2026riskreview} show that risk ranking
and review-value ranking differ when errors differ in repairability. Our
setting isolates a different source of the same gap, cost asymmetry
\emph{within} an action: wrongly accepted serious bugs are expensive, but
they share the action, the error history and largely the confidence of
the correctly accepted ones, so neither the action-level rule nor the
confidence signals we tested separated them.

\para{Lessons for practitioners} First, a per-step confidence is not a
workflow-level risk signal; the workflow's own action and that action's
error history predicted failures better for every model we evaluated.
Second, a better failure ranking should not be deployed as a review policy
without scoring it in cost: here the most error-prone action was a cheap
one, ranking by error rate spent reviews where they saved nothing, and
whether any review paid depended on the ratio of wrong-accept cost to
review cost.
Third, where review value exists it may sit in errors that no available
signal separates; a pre-registered feasibility check can establish this
before data are spent. Simple process checks also caught most of our own
design defects before data existed (\cref{app:audits}).

\para{What would make review pay} The oracle bounds the prize at 33--44\%
of cost for the Jev-style models. Capturing it needs a signal that
separates serious from minor bugs \emph{among accepted cases}, for example
evidence the workflow does not currently read, or a primitive whose
severity confidence tracks severity. Whether such a signal exists is an
open question that our data cannot answer.

\section{Conclusion}

What a typed-decision workflow \emph{decided} predicted its failures
better than how confident it was, on 500 new bugs and across five models.
That signal ranks errors by how often they happen, not by what they cost.
Under the pre-registered costs it did not make review pay, because the
expensive errors hide inside actions that look alike; whether review pays
at all turned on the ratio of wrong-accept cost to review cost; and a
case-level score aimed at the expensive errors failed a pre-registered
feasibility check, so we stopped by rule. Workflow uncertainty has to be
evaluated against the decision it is meant to improve, in that decision's
costs.

\section*{Limitations}

\para{Scope} Five models in one domain (bug triage) and one workflow. Jev
was version 1.13.0; later versions and other domains may behave
differently. The general instruct models use our prompt and readout, not a
trained decision interface, so their absolute failure rates are not a
clean comparison.

\para{Label-prior shift} Decision-conditioned error estimators fitted on
the quota-sampled development bugs can invert on a natural-mix set
(\cref{tab:targets}). The re-weighting that repairs this uses the
held-out sample's label counts (not its labels per bug or model outputs),
which is optimistic: a deployer has to estimate the mix.

\para{Labels and target} Reference labels are not ground truth: Mozilla's
sample excludes re-triaged bugs, and kernel severity is the reporter's
claim, with ``normal'' as a default. The target action is a construct we
defined, and all strategy results are conditional on it. Old public bugs
may appear in the models' training data.

\para{Reuse of the held-out set} The 500 held-out bugs were fresh for the
prediction tests only. The intervention test was pre-registered after
their prediction results were known (no cost outcome had been computed),
and the cost-sensitivity analysis is post hoc on the same bugs, so neither
is an independent confirmation.

\para{Intervention and costs} One budget and one cost table were frozen in
advance. The sensitivity analysis is exploratory and reports point
estimates; real operational costs are unknown. Reviewers are assumed
correct, and the oracle is an upper bound, not a strategy. The feasibility
check is a simulation on development data and is not confirmatory evidence
in either direction.

\para{Protocol} The cross-model study has three disclosed amendments,
including a readout bug fixed after interim results were seen. The
development protocol went through twelve dated revisions, each made
before the data it affects existed.

\section*{Ethical considerations}

All bug reports are public; texts were privacy-scrubbed, and the released
data will contain identifiers and hashes rather than report text. No
human subjects were involved. API and GPU costs were modest. AI
assistants were used for code, analysis scripts and drafting; the authors
reviewed all claims, and every reported number is regenerated from
committed results by a traceability script.
